% 05_attacks.tex — mirrors docs/07_Attack_Library.md

\section{Attack Library}\label{sec:attacks}

We ship 50 concrete attack classes (48 per-row STRIDE plus 2
cross-cutting) in our open attack library
(\texttt{attacks/}). Each class has a CVSS~v3.1 vector, a
pytest reproducibility contract, and a one-class-per-row
mapping to the STRIDE tables in \S\ref{sec:threat_model}.

\subsection{Coverage}

\begin{table}[t]
\caption{Attack library coverage. Each boundary has 6 STRIDE
letters (Spoofing through Elevation), giving 48 boundary rows;
two cross-cutting attacks (prompt injection + grammar fuzzing)
round the library out to 50 classes. The \emph{Tickets} column
lists the per-boundary ticket IDs covered; the CVSS column
lists the per-row CVSS~v3.1 vector family.}
\label{tab:coverage}
\centering
\small
\begin{tabularx}{\linewidth}{l l X X}
\toprule
Boundary & STRIDE & Tickets covered & CVSS family \\
\midrule
transport    & S/T/R/I/D/E & \texttt{A-TRN-001..009} & \texttt{AV:N/AC:L/PR:N/UI:N/S:U/...} \\
session      & S/T/R/I/D/E & \texttt{A-SES-001..009} & \texttt{AV:N/AC:H/PR:N/UI:N/S:U/...} \\
namespace    & S/T/R/I/D/E & \texttt{A-NSP-001..007} & \texttt{AV:N/AC:L/PR:L/UI:N/S:U/...} \\
tools        & S/T/R/I/D/E & \texttt{A-TOL-001..010} & \texttt{AV:N/AC:L/PR:L/UI:R/S:U/...} \\
resources    & S/T/R/I/D/E & \texttt{A-RES-001..007} & \texttt{AV:N/AC:L/PR:L/UI:N/S:U/...} \\
memory       & S/T/R/I/D/E & \texttt{A-MEM-001..008} & \texttt{AV:N/AC:L/PR:L/UI:N/S:U/...} \\
cache        & T/R/I/D/E   & \texttt{A-CCH-001..005} & \texttt{AV:N/AC:H/PR:N/UI:N/S:U/...} \\
auth         & S/T/R/I/D/E & \texttt{A-AUT-001..008} & \texttt{AV:N/AC:L/PR:N/UI:N/S:U/...} \\
prompt inj.  & n/a         & \texttt{A-TOL-001}, \texttt{A-RES-006/007} & \texttt{AV:N/AC:L/PR:L/UI:R/S:C/C:H/I:H/A:N} \\
fuzzing      & n/a         & surface-wide          & \texttt{AV:N/AC:H/PR:N/UI:N/S:U/C:H/I:H/A:H} \\
\bottomrule
\end{tabularx}
\end{table}

\subsection{Attack patterns}

We group the 48 per-row attacks into 8 reusable patterns, each
matching one boundary's failure mode:

\begin{enumerate}
  \item \textbf{\texttt{framing\_or\_replay\_inject}}
        (transport $\times$ S/T/R/I/D/E). Smuggle a payload
        marker into the framing bytes or replay surface; the
        source connector invokes \texttt{echo} with a
        marker-carrying argument set and the Evaluator checks
        whether the marker surfaces in the sink's view.
  \item \textbf{\texttt{session\_fixation}} (session $\times$
        S/T/R/I/D/E). Force a fixed session ID on the source
        connector and inspect whether the sink's session store
        returns the source's payload marker. Mirrors misuse
        case MC-2.
  \item \textbf{\texttt{tool\_shadow\_inject}} (namespace
        $\times$ S/T/R/I/D/E). Invoke a tool whose name
        shadows a built-in (\texttt{shadow\_set\_env}) and read
        a resource whose URI is shaped by the squatted name.
        Mirrors misuse case MC-1.
  \item \textbf{\texttt{result\_smuggle\_or\_inject}} (tool
        $\times$ S/T/R/I/D/E). Indirect prompt injection: the
        source connector invokes \texttt{echo} with a payload
        carrying a prompt-like instruction; the sink connector
        reads a resource whose contents include the same
        marker. Cross-boundary (tool $\leftrightarrow$ resource).
  \item \textbf{\texttt{path\_traversal\_uri}} (resource
        $\times$ S/T/R/I/D/E). Path traversal: the source
        connector reads a URI with \texttt{..} traversal and
        percent-encoded slashes (\texttt{\%2F}); the sink
        connector checks whether the traversal surfaced the
        marker. Mirrors misuse case MC-3.
  \item \textbf{\texttt{scratchpad\_or\_memory\_leak}} (memory
        $\times$ S/T/R/I/D/E). The source writes a value to a
        scratchpad keyed by name only (no
        \texttt{tenant\_id}); the sink reads the same key.
        Mirrors the confused-deputy concept.
  \item \textbf{\texttt{cache\_key\_collision}} (cache
        $\times$ T/R/I/D/E). Both connectors attempt to read
        the same cache entry keyed by
        \texttt{(tool\_name, args\_hash)}; the attack succeeds
        if the sink returns the source's payload. Mirrors
        misuse case MC-4.
  \item \textbf{\texttt{token\_replay\_or\_misuse}} (auth
        $\times$ S/T/R/I/D/E). The source connector invokes a
        tool that performs an SSRF-style fetch carrying the
        original \texttt{Authorization} header; the sink
        connector reads a resource to detect whether the token
        surfaced.
\end{enumerate}

\subsection{Cross-cutting attacks}

\paragraph{\texttt{ResultInjectionAttack} (\texttt{A-TOL-001}).}
Indirect prompt injection via tool result content. Touches both
the \textbf{tool} and \textbf{resource} boundaries. Implements
\texttt{A-TOL-001}, \texttt{A-RES-006}, and \texttt{A-RES-007}
simultaneously. CVSS~3.1:
\texttt{AV:N/AC:L/PR:L/UI:R/S:C/C:H/I:H/A:N}. This is the
direct protocol-level embodiment of the indirect prompt
injection originally identified by Greshake et al.
(2023)~\cite{greshake_2023_indirect}.

\paragraph{\texttt{GrammarFuzzAttack} (\texttt{A-FUZZ-001}).}
Grammar-aware fuzzing harness adapted from LSPFuzz (Zhu et al.
2025~\cite{zhu_2025_lspfuzz}). Generates $N$ well-formed-but-
malicious JSON-RPC envelopes with mutated parameter shapes.
LSPFuzz's campaign over $\sim$300 LSP servers is the
methodological template; ours is the MCP-specific adaptation
that respects the eight-boundary catalogue.

\subsection{Pytest reproducibility}

The library ships with a parametrized pytest suite
(\texttt{tests/test\_attack\_library.py}) that exercises every
registered id:

\begin{itemize}
  \item \texttt{test\_minimum\_attack\_count} asserts $\geq 25$
        (we ship 50).
  \item \texttt{test\_every\_boundary\_has\_at\_least\_one\_attack}
        asserts each of the 8 \texttt{Boundary} enum values has
        $\geq 1$ attack.
  \item \texttt{test\_attack\_id\_format} asserts id matches
        \texttt{A-PREFIX-NNN}.
  \item \texttt{test\_attack\_has\_cvss} asserts every class
        has a CVSS vector attribute or docstring.
  \item \texttt{test\_attack\_subclasses\_attack\_base} asserts
        every class subclasses \texttt{attacks.base.Attack}.
  \item \texttt{test\_attack\_instantiable} asserts every class
        can be instantiated.
  \item \texttt{test\_attack\_execute\_runs} runs one
        representative attack end-to-end.
  \item \texttt{test\_no\_duplicate\_ids} asserts no id
        collisions.
\end{itemize}

The full suite passes with 204 tests in $< 0.2$s, confirming
that the library is reproducible on any machine with the pinned
dependencies.

\subsection{Attack-class lifecycle}

Each attack class follows the same lifecycle:

\begin{enumerate}
  \item \textbf{Registration.} The class is registered in
        \texttt{attacks.registry.REGISTRY} keyed by its
        \texttt{id} attribute. The registry is a
        \texttt{dict[str, Type[Attack]]} consumed by the
        Scheduler.
  \item \textbf{Instantiation.} The Scheduler instantiates the
        class with the tenant pair and the per-recipe
        parameters from the manifest. The class's
        \texttt{\_\_init\_\_} validates the parameters and
        raises on type mismatches.
  \item \textbf{Execution.} The Scheduler calls
        \texttt{execute()} with the source connector, the
        sink connector, and a seeded \texttt{random.Random}.
        The method drives the cross-tenant interaction and
        returns a \texttt{list[Event]} (call events,
        leakage events, or both).
  \item \textbf{Evaluation.} The Evaluator inspects the
        events and classifies leakage based on substring
        and sha256-prefix matches against the seeded
        marker.
  \item \textbf{Logging.} The Logger appends each event to
        the run's JSONL file with the tenant pair, the
        attack id, the iteration, and the seeded marker.
\end{enumerate}

The lifecycle is uniform across all 50 classes, which makes
the framework's pluggability contract concrete: a new attack
class needs only to subclass \texttt{attacks.base.Attack},
implement \texttt{execute()}, and register itself.

\subsection{Coverage rationale}

The 48-row coverage (8 boundaries $\times$ 6 STRIDE letters)
is not arbitrary. STRIDE is the de-facto vocabulary for
protocol-level threat modelling~\cite{stride_orig}; each
letter maps onto a distinct attack pattern at each
boundary. The two cross-cutting attacks
(\texttt{ResultInjectionAttack}, \texttt{GrammarFuzzAttack})
are added because prompt injection and grammar fuzzing are
\emph{horizontal} attack patterns that cross multiple
boundaries and cannot be pinned to a single row.

The CVSS scoring per row is documented in each attack class's
\texttt{cvss} attribute; full vectors are generated by
\texttt{scripts/gen\_phase7\_implementations.py} and follow
the FIRST specification. The scores are consistent with the
adjacent literature (Chowdhury et al.'s survey
\cite{chowdhury_2024_defenses}; Wang et al.'s ToolCommander
\cite{wang_2024_toolcommander}) and serve as a triage tool
for defenders, not as a precise impact prediction.

\subsection{Comparison with adjacent attack libraries}

The closest prior attack libraries are:

\begin{itemize}
  \item \textbf{Chowdhury et al.'s LLM-attack survey}
        \cite{chowdhury_2024_defenses}: catalogue of LLM-level
        attacks; no MCP-specific surface.
  \item \textbf{An et al.'s FlowGuard}
        \cite{an_2026_flowguard}: detection traces, not
        attack implementations.
  \item \textbf{Zhu et al.'s LSPFuzz}
        \cite{zhu_2025_lspfuzz}: grammar-aware fuzzing on
        LSP; not directly portable to MCP.
  \item \textbf{Liu et al.'s MCPEvol-Bench}
        \cite{liu_2026_mcpevol}: dynamic-evolution benchmark;
        no attack implementations.
\end{itemize}

Our library is the first to ship MCP-specific attack
implementations with pytest reproducibility, CVSS scoring,
and uniform per-(boundary, STRIDE) coverage. The combination
is the unit of novelty: prior work covers some of these
dimensions in isolation, none in composition.